% META: {"id":"1SPE-EXPO-ME-002","chapitre":"1SPE-EXPONENTIELLE","type_objet":"methode","capacites_codes":["C3"],"status":"approved","origin":"generated","status_history":["generated","audited","accepted","approved"],"release_acceptance":"RELEASE_OWNER_BATCH_ACCEPTANCE","acceptance_closure_digest":"sha256:9b3ccf9a81c5520fbb7e03b7b2d3e7bf2057a02834b6d908bf3be5f49f3b553f","acceptance_receipt_digest":"sha256:4fad86f0282e0fa4e0419cca1ad6379ae7af5a280c04a4a90880f90a1c044242"}
\begin{fichemethode}{M2}{Comparer des valeurs de la fonction exponentielle}
\quandUtiliser{J'utilise cette méthode pour comparer deux exponentielles ou situer une valeur sur la courbe, sans calculatrice.}
\pasApas{\begin{enumerate}
  \item Je compare les deux exposants.
  \item J'utilise la stricte croissance de la fonction exponentielle : l'ordre est conservé.
  \item Je contrôle le signe : toute valeur exponentielle est strictement positive.
\end{enumerate}}
\exempleRedige{Comparer $\mathrm{e}^{2{,}5}$ et $\mathrm{e}^{\sqrt7}$.
Comme $2{,}5^2=6{,}25<7$, on a $2{,}5<\sqrt7$. La fonction exponentielle
étant strictement croissante,
\[\mathrm{e}^{2{,}5}<\mathrm{e}^{\sqrt7}.\]}
\pieges{\begin{itemize}
  \item Comparer des valeurs approchées alors que les exposants se comparent exactement.
  \item Inverser l'ordre : la fonction exponentielle est croissante.
  \item Croire qu'un exposant négatif donne une valeur négative.
\end{itemize}}
\verifier{Je peux contrôler sur la courbe que l'ordonnée augmente avec l'abscisse.}
\sEntrainer{\refExos{M2}}
\end{fichemethode}
